\subsection{Integration with respect to a measure with compact support}

Let \(\mu\) be a measure on X whose support \(S = \operatorname{Supp}(\mu)\) is compact; the open set \(X - S\) is negligible (§2, No.~2, Prop. 5). For every function \(\mathbf{f}\) with values in a vector space \(\mathbf{F}\) or in \(\overline{\mathbf{R}}\) , the functions \(\mathbf{f}\) and \(\mathbf{f}\varphi_{\mathrm{S}}\) are therefore equivalent (§2, No.~4); for \(\mathbf{f}\) to be \(\mu\) -integrable (when \(\mathbf{F}\) is a Banach space), it is therefore necessary and sufficient that \(\mathbf{f}\varphi_{\mathrm{S}}\) be so, in which case (No.~1)

\[
\int \mathbf {f} d \mu = \int \mathbf {f} \varphi_ {\mathrm{S}} d \mu .\tag{21}
\]

If, moreover, \(\mathbf{f}\) is bounded on \(S\) , it follows from (20) that

\[
\left| \int \mathbf {f} d \mu \right| \leqslant \| \mu \| \cdot \sup _ {x \in S} | \mathbf {f} (x) |.\tag{22}
\]

In particular, if \(\mathbf{f}\) is continuous on \(X\) then \(\mathbf{f}\) is \(\mu\) -integrable, since \(\mathbf{f}h \in \mathcal{K}(X;F)\) for every function \(h \in \mathcal{K}(X;R)\) equal to 1 on \(S\) (Ch. III, §1, No.~2, Lemma 1). More precisely:

PROPOSITION 14. --- Let X be a locally compact space, F a Banach space not reduced to 0; equip the space \(\mathcal{C}(\mathrm{X};\mathrm{F})\) of all continuous mappings of X into F with the topology of compact convergence. For a measure \(\mu\) on X to be such that the linear mapping \(\mathbf{f} \mapsto \int \mathbf{f} d\mu\) of \(\mathcal{K}(\mathrm{X};\mathrm{F})\) into F is extendible to a continuous linear mapping of \(\mathcal{C}(\mathrm{X};\mathrm{F})\) into F, it is necessary and sufficient that \(\operatorname{Supp}(\mu)\) be compact; such an extension is unique and coincides with the integral defined in No.~1.

We have just seen that if \(\mu\) has compact support, then the integral \(\int \mathbf{f}d\mu\) is defined for every function \(\mathbf{f}\in \mathcal{C}(\mathrm{X};\mathrm{F})\) and that the mapping \(\mathbf{f}\mapsto \int \mathbf{f}d\mu\) of \(\mathcal{C}(\mathrm{X};\mathrm{F})\) into \(\mathrm{F}\) is continuous for the topology of compact convergence. Conversely, suppose that \(\mathbf{f}\mapsto \int \mathbf{f}d\mu\) is continuous in \(\mathcal{K}(\mathrm{X};\mathrm{F})\) for the topology of compact convergence. Then, there is a compact set \(K\subset X\) and a number \(a > 0\) such that \(|\mu (\mathbf{f})|\leqslant a\cdot \sup_{x\in K}|\mathbf{f}(x)|\) for every function \(\mathbf{f}\in \mathcal{K}(\mathrm{X};\mathrm{F})\) ; in particular, if the support of \(\mathbf{g}\in \mathcal{K}(\mathrm{X};\mathrm{F})\) does not intersect \(K\) , then \(\mu (\mathbf{g}) = 0\) . Taking \(\mathbf{g} = h\mathbf{a}\) , where \(\mathbf{a}\neq 0\) is a vector in \(\mathrm{F}\) and \(h\in \mathcal{K}(\mathrm{X};\mathbf{C})\) , we see that \(\mu (h) = 0\) for every function \(h\in \mathcal{K}(\mathrm{X};\mathbf{C})\) whose support does not intersect \(K\) , which proves that \(\operatorname {Supp}(\mu)\subset K\) . Finally, the uniqueness of the extension follows from the fact that \(\mathcal{K}(\mathrm{X};\mathrm{F})\) is dense in \(\mathcal{C}(\mathrm{X};\mathrm{F})\) for the topology of compact convergence (Ch. III, §1, No.~2, Prop. 4).

Prop. 14 permits identifying a measure on X with compact support with its continuous extension to \(\mathcal{C}(\mathrm{X};\mathbf{C})\) . The set of measures on X with compact support may therefore be identified with the dual \(\mathcal{C}'(\mathrm{X};\mathbf{C})\) of the Hausdorff locally convex space \(\mathcal{C}(\mathrm{X};\mathbf{C})\) . Recall that \(\mathcal{C}(\mathrm{X};\mathbf{C})\) is complete (GT, X, §1, No.~6, Cor. 3 of Th. 2), but it is not necessarily barreled (Exer. 17). However, if X is countable at infinity, hence is the union of an increasing sequence of compact sets \(\mathrm{K}_n\) such that \(\mathrm{K}_n \subset \overset{\circ}{\mathrm{K}}_{n+1}\) , then the topology of \(\mathcal{C}(\mathrm{X};\mathbf{C})\) can be defined by the countable family of semi-norms \(p_n(f) = \sup_{x \in \mathrm{K}_n} |f(x)|\) , therefore \(\mathcal{C}(\mathrm{X};\mathbf{C})\) is a Fréchet space in this case. Consequently, for every covering \(\mathfrak{S}\) of \(\mathcal{C}(\mathrm{X};\mathbf{C})\) by bounded sets, the space \(\mathcal{C}'(\mathrm{X};\mathbf{C})\) is then quasi-complete for the \(\mathfrak{S}\) -topology (TVS, III, §4, No.~2, Cor. 4 of Th. 1).

We shall consider above all on \(\mathcal{C}^{\prime}(\mathrm{X};\mathbf{C})\) the topology of compact convergence (the topology of uniform convergence on the compact subsets of \(\mathcal{C}(\mathrm{X};\mathbf{C})\) ). Recall that the relatively compact subsets H of \(\mathcal{C}(\mathrm{X};\mathbf{C})\) are characterized by the following properties (GT, X, §2, No.~5, Cor. 3 of Th. 2):

\(1^{\circ}\) H is equicontinuous;

\(2^{\circ}\) for every \(x \in \mathrm{X}\) , the set \(\mathrm{H}(x)\) of the \(f(x)\) , where \(f\) runs over \(\mathrm{H}\) , is bounded in \(\mathbf{C}\) .

PROPOSITION 15. --- Let X be a locally compact space and, for every \(x \in X\) , let \(\varepsilon_{x}\) be the Dirac measure at the point x. The mapping \(x \mapsto \varepsilon_{x}\) of X into \(\mathcal{C}'(X; \mathbf{C})\) is continuous for the topology of compact convergence on \(\mathcal{C}'(X; \mathbf{C})\) .

Consider a neighborhood of \(\varepsilon_{x_0}\) in \(\mathcal{C}'(\mathrm{X};\mathbf{C})\) for this topology, which we can suppose to be defined by taking a number \(\delta > 0\) , a compact subset \(\mathrm{H}\) of \(\mathcal{C}(\mathrm{X};\mathbf{C})\) , and considering the set of measures \(\mu\) on \(\mathrm{X}\) with compact support such that \(|\mu(f) - \varepsilon_{x_0}(f)| \leqslant \delta\) for every function \(f \in \mathrm{H}\) . Since \(\mathrm{H}\) is equicontinuous, there exists a neighborhood \(\mathrm{U}\) of \(x_0\) in \(\mathrm{X}\) such that the relation \(f \in \mathrm{H}\) implies \(|f(x) - f(x_0)| \leqslant \delta\) for all \(x \in \mathrm{U}\) , which may also be written \(|\varepsilon_x(f) - \varepsilon_{x_0}(f)| \leqslant \delta\) and proves the proposition. (*)

PROPOSITION 16. --- Let K be a compact subset of X, L the vector space of measures \(\mu\) on X with support contained in K. On L, the topologies induced by the topology \(\mathcal{T}\) of compact convergence on \(\mathcal{C}'(X; \mathbf{C})\) and the topology \(\mathcal{T}'\) of strictly compact convergence on \(\mathcal{M}(X; \mathbf{C})\) (Ch. III, §1, No.~10) coincide.

It is clear that on L, the topology induced by T is finer than the topology induced by \(T'\) . Conversely, let H be a compact subset of \(\mathcal{C}(\mathrm{X};\mathbf{C})\) , h a function in \(\mathcal{K}(\mathrm{X};\mathbf{C})\) equal to 1 on K. It is clear that the set \(H'\) of functions fh, where f runs over H, is strictly compact in \(\mathcal{K}(\mathrm{X};\mathbf{C})\) , and, for every measure \(\mu\in L\) , \(\mu(f)=\mu(fh)\) for every function \(f\in H\) , whence the conclusion.

COROLLARY 1. --- For every compact subset K of X and every number a \textgreater{} 0, the set B of measures \(\mu\) on X such that \(\operatorname{Supp}(\mu) \subset K\) and \(\|\mu\| \leqslant a\) is an equicontinuous subset of \(\mathcal{C}'(X; \mathbf{C})\) that is compact for the topology T of compact convergence.

For, let H be a subset of \(\mathcal{C}(\mathrm{X};\mathbf{C})\) consisting of functions that are uniformly bounded on K; there exists a number c\textgreater0 such that \(|\mu(f)|\leqslant c\cdot\|\mu\|\leqslant ac\) for every function \(f\in H\) and every measure \(\mu\in B\) , by virtue of (22); therefore \(B\subset acH^{\circ}\) in the dual \(\mathcal{C}'(\mathrm{X};\mathbf{C})\) of \(\mathcal{C}(\mathrm{X};\mathbf{C})\) , which proves the equicontinuity of B; the fact that B is compact for T follows from the fact that, on B, T and the vague topology induce the same topology (Prop. 16 and Ch. III, §1, No.~10, Prop. 17) and the fact that B is vaguely compact (Ch. III, §1, No.~9, Cor. 2 of Prop. 15 and §2, No.~2, Prop. 6).

COROLLARY 2. --- Every measure with compact support (resp. every positive measure with compact support) \(\mu\) is in the closure in \(\mathcal{C}'(\mathrm{X};\mathbf{C})\) , for the topology \(\mathcal{T}\) of compact convergence, of the set of measures (resp.

positive measures) whose support is finite and contained in \(\operatorname{Supp}(\mu)\) and whose norm is equal to \(\| \mu \|\) .

For, on the set B of measures \(\nu\) such that \(\operatorname{Supp}(\nu) \subset \operatorname{Supp}(\mu)\) and \(\|\nu\| \leqslant \|\mu\|\) , the topology induced by the vague topology is identical to the topology induced by T, and the corollary therefore follows from Ch. III, §2, No.~4, Cors. 2 and 3 of Th. 1.

